\Article{18}
Wenden wir die vorangegangene Untersuchung auf die Auswertung des zweiten Terms des Ausdrucks $\Omega$ in Artikel 3 an, so ist der zweite, seit Artikel 6 mit $S$ bezeichnete Raumbereich mit dem ersten identisch. Was also in Artikel 16 $\sigma$, $\tau$, $\tau'$ war, ist hier $s$, $t$, $0$, wenn $t$ die ganze Oberfläche des von der Flüssigkeit erfüllten Raumbereichs $s$ bezeichnet. Enthält dieser Raumbereich weder Teile von merklicher Ausdehnung und unmerklicher Dicke noch Zwischenräume (Spalten) solcher Art, so ergibt sich für den zweiten Teil des Ausdrucks $\Omega$
\[
=\frac12\pi cc(s\psi0-t\theta0)
\]

Es gibt also zwei Ausnahmen.

1) Enthält der Raumbereich $s$ einen Teil von unmerklicher Dicke, so weist dessen Oberfläche zwei Teile auf, die sich in ihrer Größe nicht merklich unterscheiden. Bezeichnen wir einen von beiden mit $t'$ und die Dicke des Raumbereichs bei einem beliebigen Element $\mathrm{d}t'$ allgemein mit $\rho$, so kommt zum vorigen Ausdruck der Term
\[
\pi cc\int\theta'\rho.\mathrm{d}t'
\]
hinzu.

2) Enthält der Raumbereich $s$ einen Hohlraum von unmerklicher Dicke, so kommt ein ähnlicher Term hinzu, nämlich $\pi cc\int\theta'\rho.\mathrm{d}t''$. Dabei bezeichnet $t''$ einen der beiden an den Spalt grenzenden Teile der Oberfläche $t$ und $\rho$ allgemein die Dicke des Spalts an einem beliebigen Punkt.

Bei der Auswertung des dritten Terms des Ausdrucks $\Omega$ behalten wir das Zeichen $S$ für den vom Gefäß eingenommenen Raumbereich bei. An die Stelle des Funktionszeichens $f$ muss jedoch $F$ \SourceLeaf{58}für die Anziehungskraft der Moleküle des Gefäßes treten. Ebenso verwenden wir anstelle der mit $\varphi$, $\psi$, $\theta$, $\theta'$ bezeichneten Funktionen andere, die wir mit $\Phi$, $\Psi$, $\Theta$, $\Theta'$ bezeichnen und von denen wir annehmen, dass sie ebenso von $F$ abhängen wie jene von $f$. Was in der allgemeinen Untersuchung $\sigma$, $\tau'$ war, ist hier offenbar $0$. Für $\tau$ verwenden wir hier einfach den Buchstaben $T$, um von der Oberfläche des Raumbereichs $S$ nur den Teil zu bezeichnen, der an die Flüssigkeit grenzt. Damit ergibt sich für den dritten Teil des Ausdrucks $\Omega$ im Allgemeinen
\[
=\pi cCT\Theta0
\]
wobei auch hier zwei Fälle auszunehmen sind.

3) Hat die Flüssigkeit an einem merklich großen Teil $T'$ der Oberfläche $T$ eine unmerkliche Dicke, die allgemein durch $\rho$ bezeichnet werde, so kommt der Term
\[
-\pi cC\int\Theta'\rho.\mathrm{d}T'
\]
hinzu.

4) Weist die Oberfläche des Gefäßes neben dem an die Flüssigkeit grenzenden Teil $T$ einen weiteren Teil $T''$ auf, der zwar in einem gewissen, aber unmerklichen Abstand von der Flüssigkeit liegt, so kommt, wenn $\rho$ allgemein diesen Abstand für einen beliebigen Punkt bezeichnet, der Term
\[
+\pi cC\int\Theta'\rho.\mathrm{d}T''
\]
hinzu. Es wäre überflüssig, bei der ersten Ausnahme, soweit sie in der dritten nicht enthalten ist, oder bei der zweiten oder vierten zu verweilen. Ein Gleichgewicht der Flüssigkeit könnte zwar in gewissen hierher gehörenden, doch höchst besonderen Fällen vorkommen, aber ein solches Gleichgewicht könnte weder stabil noch Versuchen zugänglich sein. Dagegen ist der erste Ausnahmefall, soweit er im dritten enthalten ist, durchaus wesentlich für die Theorie. Er wird hier dennoch eine Weile beiseitegelassen, damit wir die Bedingungen für ein Gleichgewicht untersuchen können, soweit es ohne eine am Gefäß haftende unmerkliche Flüssigkeitshaut bestehen kann.

Lassen wir also alle diese Ausnahmen beiseite, so lautet der Ausdruck, dessen Wert im Gleichgewichtszustand ein Maximum sein muss,
\[
-gc\int z\mathrm{d}s+\frac12ccs\psi0-\frac12\pi cct\theta0+\pi cCT\Theta0
\]
und da das Volumen $s$ bei allen Veränderungen, die die Gestalt der Flüssigkeit erfahren kann, unverändert bleibt, muss der folgende Ausdruck
\[
\int z\mathrm{d}s+\frac{\pi c\theta0}{2g}.t-\frac{\pi C\Theta0}{g}.T
\]
im Gleichgewichtszustand ein \textit{Minimum} sein.

\SourceLeaf{59}Wir haben bereits darauf hingewiesen, dass $\frac{c\theta0}{g}$ einen zweidimensionalen Bereich darstellt. Dasselbe gilt für $\frac{C\Theta0}{g}$. Setzen wir also
\[
\frac{\pi c\theta0}{2g}=\alpha\alpha,\qquad\frac{\pi C\Theta0}{2g}=\beta\beta
\]
so sind $\alpha$, $\beta$ konstante Längen, die vom Verhältnis der Schwerkraft zur Stärke der Kräfte abhängen, die die Teile der Flüssigkeit aufeinander und die Moleküle des Gefäßes auf sie ausüben. Bezeichnen wir ferner den freien Teil der Flüssigkeitsoberfläche, also den Teil, der nicht an das Gefäß grenzt, mit $U$, sodass $t=T+U$ gilt, so muss der folgende, von nun an mit $W$ bezeichnete Ausdruck im Gleichgewichtszustand ein Minimum sein
\[
\int z\mathrm{d}s+(\alpha\alpha-2\beta\beta)T+\alpha\alpha U
\]

\Article{19}
Bevor wir die Folgerungen aus diesem wichtigen Satz allgemein und vollständig herleiten, lohnt es sich zu zeigen, wie leicht sich die Haupterscheinung in Kapillarröhren daraus ergibt.

Betrachten wir eine Flüssigkeit im Gleichgewicht in einem Gefäß mit zwei Schenkeln, wobei ein Teil der freien Oberfläche der Flüssigkeit im ersten Schenkel und ein anderer Teil im zweiten liegt. Wir nehmen an, dass die Gefäßwände dort, wo sie an diese Teile grenzen, senkrecht sind. Es sei $a$ der Flächeninhalt des inneren horizontalen Querschnitts des ersten Schenkels (oder genauer der horizontalen Projektion der freien Flüssigkeitsoberfläche im ersten Schenkel), $b$ sein Umfang und schließlich $ah$ das Volumen der Flüssigkeit in diesem Schenkel, wenn die senkrechte Wand nach unten bis zu der Ebene fortgesetzt wird, von der aus die Abstände $z$ gemessen werden. Gleichbedeutend damit ist $h$ die mittlere Höhe der Flüssigkeit über dieser Ebene. Die entsprechenden Größen für den zweiten Schenkel seien mit $a'$, $b'$, $h'$ bezeichnet. Denken wir uns, dass der Zustand der Flüssigkeit eine unendlich kleine Änderung erfährt, bei der beide Teile der freien Oberfläche ihre Gestalt behalten, so ist die Variation des ersten Teils des Ausdrucks $W$, nämlich des Integrals $\int z\mathrm{d}s$, offenbar
\[
=ah\mathrm{d}h+a'h'\mathrm{d}h
\]
die Variation von $T$ dagegen
\[
=b\mathrm{d}h+b'\mathrm{d}h'
\]
und schließlich gilt nach Voraussetzung $\mathrm{d}U=0$. Daraus folgt
\[
\mathrm{d}W=ah\mathrm{d}h+a'h'\mathrm{d}h'-(2\beta\beta-\alpha\alpha)(b\mathrm{d}h+b'\mathrm{d}h')
\]
\SourceLeaf{60}Da ferner das Gesamtvolumen der Flüssigkeit unverändert bleibt, gilt
\[
a\mathrm{d}h+a'\mathrm{d}h'=0
\]
und daher
\[
\mathrm{d}W=\mathrm{d}h[a(h-h')-(2\beta\beta-\alpha\alpha)(b-\frac{ab'}{a'})]
\]
Die Bedingung, dass $W$ im Gleichgewichtszustand ein Minimum ist, führt also zu der Gleichung, die die Haupterscheinung in Kapillarröhren erfasst
\[
h-h'=(2\beta\beta-\alpha\alpha)(\frac{b}{a}-\frac{b'}{a'})
\]
und man sieht unmittelbar, dass dieser Gleichung tatsächlich ein \textit{Minimum} von $W$ entspricht, denn $\frac{\mathrm{d}\mathrm{d}W}{\mathrm{d}h^2}$ hat den Wert $=a+\frac{aa}{a'}$, der seiner Natur nach positiv ist.

Der zweite Schenkel heißt weiter als der erste, wenn $\frac{a'}{b'}$ größer als $\frac{a}{b}$ ist. Die Flüssigkeit ist also im engeren Schenkel stärker abgesenkt oder angehoben als im weiteren, je nachdem, ob das Quadrat $\beta\beta$ kleiner oder größer als $\frac12\alpha\alpha$ ist. Sollte sich $\beta\beta=\frac12\alpha\alpha$ ergeben, so wäre die Höhe in beiden Schenkeln gleich. Ist der zweite Schenkel so weit, dass $\frac{b'}{a'}$ gegenüber $\frac{b}{a}$ vernachlässigt werden kann, so gilt näherungsweise
\[
h-h'=(2\beta\beta-\alpha\alpha)\frac{b}{a}
\]
In zylindrischen Kapillarröhren ist demnach die Absenkung oder der Anstieg der Flüssigkeit umgekehrt proportional zum Durchmesser der Röhre. All dies stimmt sowohl mit der Erfahrung als auch mit dem überein, was der berühmte \textsc{Laplace} theoretisch zu begründen versucht hat.

Hat das Gefäß mehrere senkrechte Schenkel, die miteinander in Verbindung stehen, so mögen $a''$, $b''$, $h''$ für den dritten und $a'''$, $b'''$, $h'''$ für den vierten usw. dasselbe bezeichnen wie $a$, $b$, $h$ für den ersten. Dann gilt auch
\[
\begin{gathered}
h-h''=(2\beta\beta-\alpha\alpha)(\frac{b}{a}-\frac{b''}{a'})\\
h-h'''=(2\beta\beta-\alpha\alpha)(\frac{b}{a}-\frac{b'''}{a'''})
\end{gathered}
\]
Gefälliger lassen sich diese Gleichungen folgendermaßen schreiben
\[
\begin{aligned}
h-(2\beta\beta-\alpha\alpha)\frac{b}{a}&=h'-(2\beta\beta-\alpha\alpha)\frac{b'}{a'}\\
&=h''-(2\beta\beta-\alpha\alpha)\frac{b''}{a''}=h'''-(2\beta\beta-\alpha\alpha)\frac{b'''}{a'''}\text{etc.}
\end{aligned}
\]
Da die horizontale Ebene, von der aus die Höhen gemessen werden, beliebig ist, gilt, wenn wir sie so wählen, dass
\[
h=(2\beta\beta-\alpha\alpha)\frac{b}{a}
\]
\SourceLeaf{61}offenbar auch in den übrigen Schenkeln
\[
\begin{gathered}
h'=(2\beta\beta-\alpha\alpha)\frac{b'}{a'},\qquad h''=(2\beta\beta-\alpha\alpha)\frac{b''}{a''},\\
h'''=(2\beta\beta-\alpha\alpha)\frac{b'''}{a'''}\quad\text{etc.}
\end{gathered}
\]
Diese Ebene, deren Begriff wir weiter unten allgemeiner bestimmen werden, kann normale Horizontalebene (plan de niveau) heißen. Denken wir uns die senkrechten Wände der einzelnen Schenkel nötigenfalls bis zu dieser Ebene verlängert, so geben $ah$, $a'h'$, $a''h''$ usw. für $2\beta\beta>\alpha\alpha$ die Mengen der in den einzelnen Schenkeln über diese Ebene gehobenen Flüssigkeit an, für $2\beta\beta<\alpha\alpha$ dagegen die in den einzelnen Schenkeln unterhalb dieser Ebene fehlenden Flüssigkeitsmengen. Diese Mengen sind also gleich den Produkten aus der konstanten Fläche $2\beta\beta-\alpha\alpha$ und den Umfängen $b$, $b'$, $b''$, $b'''$ usw.

